\section{Prior work on ML benchmarks for distribution shifts}\label{sec:app_related}

In this section, we discuss existing ML distribution shift benchmarks in more detail, categorizing them by how they induce their respective distribution shifts.
We focus here on work that has appeared in ML conferences and journals; we discuss related work from other research communities in \refsec{other_datasets} and \refapp{app_datasets}.
We also restrict our attention to publicly-available datasets.
While others have studied some proprietary datasets with realistic distribution shifts, such as the StreetView StoreFronts dataset \citep{hendrycks2020many} or diabetic retinopathy datasets \citep{damour2020underspecification}, these datasets are not publicly available due to privacy and other commercial reasons.

\paragraph{Distribution shifts from transformations.} Some of the most widely-adopted benchmarks induce distribution shifts by synthetically transforming the data.
Examples include rotated and translated versions of MNIST and CIFAR \citep{worrall2017harmonic,gulrajani2020search};
surface variations such as texture, color, and corruptions like blur in Colored MNIST \citep{gulrajani2020search}, Stylized ImageNet \citep{geirhos2018imagenet}, ImageNet-C \citep{hendrycks2019benchmarking},
and similar ImageNet variants  \citep{geirhos2018generalisation};
and datasets that crop out objects and replace their backgrounds, as in the Backgrounds Challenge \citep{xiao2020noise} and other similar datasets \citep{sagawa2020group, koh2020bottleneck}.
Benchmarks for adversarial robustness also fall in this category of distribution shifts from transformations \citep{goodfellow2015explaining,croce2020robustbench}.
Though adversarial robustness is not a focus of this work,
we note that recent work on temporal perturbations with the ImageNet-Vid-Robust and YTBB-Robust datasets \citep{shankar2019image}
represents a different form of distribution shift that also impacts real-world applications.
Outside of visual object recognition, other work has used synthetic datasets and transformations to explore compositional generalization, e.g., SCAN \citep{lake2018generalization}.
We discuss this more in \refsec{other_datasets}.

\paragraph{Synthetic-to-real transfers.}
Fully synthetic datasets such as SYNTHIA \citep{ros2016synthia} and StreetHazards \citep{dan2020scaling} have been adopted for out-of-distribution detection as well as domain adaptation and generalization, e.g., by testing robustness to transformations in the seasons, weather, time, or architectural style \citep{hoffman2018cycada,volpi2018generalizing}.
While the data is synthetic, it can still look realistic if a high-fidelity simulator is used.
In particular, synthetic benchmarks that study transfers from synthetic to real data \citep{ganin2015domain,richter2016playing,peng2018visda} can be important tools for tackling real-world problems: even though the data is synthesized and by definition, not real, the synthetic-to-real distribution shift can still be realistic in contexts where real data is much harder to acquire than synthetic data \citep{bellemare2020autonomous}.
In this work, we do not study these types of synthetic distribution shifts; instead, we focus on distribution shifts that occur in the wild between real data distributions.

\paragraph{Distribution shifts from constrained splits.} Other benchmarks do not rely on transformations but instead split the data in a way that induces particular distribution shifts.
These benchmarks have realistic data, e.g., the data points are derived from real-world photos,
but they do not necessarily reflect distribution shifts that would arise in the wild.
For example, BREEDS \citep{santurkar2020breeds} and a related dataset \citep{hendrycks2019benchmarking} test generalization to unseen ImageNet subclasses by holding out subclasses specified by several controllable parameters;
similarly, NICO \citep{he2020towards} considers subclasses that are defined by their context, such as dogs at home versus dogs on the beach;
DeepFashion-Remixed \citep{hendrycks2020many} constrains the training set to include only photos from a single camera viewpoint and tests generalization to unseen camera viewpoints;
BDD-Anomaly \citep{dan2020scaling} uses a driving dataset but with all motorcycles, trains, and bicycles removed from the training set only;
and ObjectNet \citep{barbu2019objectnet} comprises images taken from a few pre-specified viewpoints, allowing for systematic evaluation for robustness to camera angle changes but deviating from natural camera angles.

\paragraph{Distribution shifts across datasets.} A well-studied special case of the above category is the class of distribution shifts obtained by combining several disparate datasets \citep{torralba2011unbiased}, training on one or more of them and then testing on the remaining datasets.
A recent influential example is the ImageNetV2 dataset \citep{recht2019doimagenet}, which was constructed to be similar to the original ImageNet dataset.
Unlike ImageNetV2, however, many of these distribution shifts were constructed to be more drastic than might arise in the wild.
For example, standard domain adaptation benchmarks include training on MNIST but testing on SVHN street signs \citep{lecun1998mnist,netzer2011reading,tzeng2017domain,hoffman2018cycada},
as well as transfers across datasets containing different renditions (e.g., photos, clipart, sketches) in DomainNet \citep{peng2019moment} and the Office-Home dataset \citep{venkateswara2017deep}.

The main difference between domain adaptation and domain generalization is that in the latter, we do not assume access to unlabeled data from the test distribution.
This makes it straightforward to use domain adaptation benchmarks for domain generalization, e.g., in DomainBed \citep{gulrajani2020search}; we focus on domain generalization in this work,
but further discuss unsupervised domain adaptation in \refsec{other_problems}.
Other similar benchmarks that have been proposed for domain generalization include VLCS \citep{fang2013unbiased}, which tests generalization across similar visual object recognition datasets; PACS \citep{li2017deeper}, which (like DomainNet) tests generalization across datasets with different renditions; and ImageNet-R \citep{hendrycks2020many} and ImageNet-Sketch \citep{wang2019learning}, which also test generalization across different renditions by collecting separate datasets from Flickr and Google Image queries.
